\section{Introdução}
\label{sec:introducao}

Esta pesquisa investiga a relação entre exposição ocupacional à inteligência artificial e transições no mercado de trabalho brasileiro. O interesse está nas mudanças de ocupação, na saída do emprego e na passagem para a informalidade, observadas em um painel construído com a PNAD Contínua. O lançamento público do ChatGPT é utilizado como referência temporal para organizar a comparação entre períodos.

A pergunta central é se trabalhadores em ocupações com diferentes níveis de exposição à inteligência artificial apresentam padrões distintos de transição após esse marco. A discussão sobre automação e criação de tarefas oferece uma referência para formular o problema, conforme \citeonline{acemoglu2019}. Entretanto, a exposição ocupacional não comprova adoção efetiva da tecnologia, e a comparação temporal, isoladamente, não permite atribuir causalidade aos resultados.

Este capítulo integra um esqueleto provisório preparado para o exercício ia-6.1. Seu conteúdo será revisado pelo autor durante o desenvolvimento da dissertação, preservando a distinção entre pergunta de pesquisa, procedimentos de análise e conclusões sustentadas por evidências.